% Options for packages loaded elsewhere
\PassOptionsToPackage{unicode}{hyperref}
\PassOptionsToPackage{hyphens}{url}
\PassOptionsToPackage{dvipsnames,svgnames,x11names}{xcolor}
\documentclass[
  11pt,
]{article}
\usepackage{xcolor}
\usepackage[a4paper,margin=2.2cm]{geometry}
\usepackage{amsmath,amssymb}
\setcounter{secnumdepth}{-\maxdimen} % remove section numbering
\usepackage{iftex}
\ifPDFTeX
  \usepackage[T1]{fontenc}
  \usepackage[utf8]{inputenc}
  \usepackage{textcomp} % provide euro and other symbols
\else % if luatex or xetex
  \usepackage{unicode-math} % this also loads fontspec
  \defaultfontfeatures{Scale=MatchLowercase}
  \defaultfontfeatures[\rmfamily]{Ligatures=TeX,Scale=1}
\fi
\usepackage{lmodern}
\ifPDFTeX\else
  % xetex/luatex font selection
\fi
% Use upquote if available, for straight quotes in verbatim environments
\IfFileExists{upquote.sty}{\usepackage{upquote}}{}
\IfFileExists{microtype.sty}{% use microtype if available
  \usepackage[]{microtype}
  \UseMicrotypeSet[protrusion]{basicmath} % disable protrusion for tt fonts
}{}
\makeatletter
\@ifundefined{KOMAClassName}{% if non-KOMA class
  \IfFileExists{parskip.sty}{%
    \usepackage{parskip}
  }{% else
    \setlength{\parindent}{0pt}
    \setlength{\parskip}{6pt plus 2pt minus 1pt}}
}{% if KOMA class
  \KOMAoptions{parskip=half}}
\makeatother
\usepackage{longtable,booktabs,array}
\newcounter{none} % for unnumbered tables
\usepackage{calc} % for calculating minipage widths
% Correct order of tables after \paragraph or \subparagraph
\usepackage{etoolbox}
\makeatletter
\patchcmd\longtable{\par}{\if@noskipsec\mbox{}\fi\par}{}{}
\makeatother
% Allow footnotes in longtable head/foot
\IfFileExists{footnotehyper.sty}{\usepackage{footnotehyper}}{\usepackage{footnote}}
\makesavenoteenv{longtable}
\usepackage{graphicx}
\makeatletter
\newsavebox\pandoc@box
\newcommand*\pandocbounded[1]{% scales image to fit in text height/width
  \sbox\pandoc@box{#1}%
  \Gscale@div\@tempa{\textheight}{\dimexpr\ht\pandoc@box+\dp\pandoc@box\relax}%
  \Gscale@div\@tempb{\linewidth}{\wd\pandoc@box}%
  \ifdim\@tempb\p@<\@tempa\p@\let\@tempa\@tempb\fi% select the smaller of both
  \ifdim\@tempa\p@<\p@\scalebox{\@tempa}{\usebox\pandoc@box}%
  \else\usebox{\pandoc@box}%
  \fi%
}
% Set default figure placement to htbp
\def\fps@figure{htbp}
\makeatother
% definitions for citeproc citations
\NewDocumentCommand\citeproctext{}{}
\NewDocumentCommand\citeproc{mm}{%
  \begingroup\def\citeproctext{#2}\cite{#1}\endgroup}
\makeatletter
 % allow citations to break across lines
 \let\@cite@ofmt\@firstofone
 % avoid brackets around text for \cite:
 \def\@biblabel#1{}
 \def\@cite#1#2{{#1\if@tempswa , #2\fi}}
\makeatother
\newlength{\cslhangindent}
\setlength{\cslhangindent}{1.5em}
\newlength{\csllabelwidth}
\setlength{\csllabelwidth}{3em}
\newenvironment{CSLReferences}[2] % #1 hanging-indent, #2 entry-spacing
 {\begin{list}{}{%
  \setlength{\itemindent}{0pt}
  \setlength{\leftmargin}{0pt}
  \setlength{\parsep}{0pt}
  % turn on hanging indent if param 1 is 1
  \ifodd #1
   \setlength{\leftmargin}{\cslhangindent}
   \setlength{\itemindent}{-1\cslhangindent}
  \fi
  % set entry spacing
  \setlength{\itemsep}{#2\baselineskip}}}
 {\end{list}}
\usepackage{calc}
\newcommand{\CSLBlock}[1]{\hfill\break\parbox[t]{\linewidth}{\strut\ignorespaces#1\strut}}
\newcommand{\CSLLeftMargin}[1]{\parbox[t]{\csllabelwidth}{\strut#1\strut}}
\newcommand{\CSLRightInline}[1]{\parbox[t]{\linewidth - \csllabelwidth}{\strut#1\strut}}
\newcommand{\CSLIndent}[1]{\hspace{\cslhangindent}#1}
\setlength{\emergencystretch}{3em} % prevent overfull lines
\providecommand{\tightlist}{%
  \setlength{\itemsep}{0pt}\setlength{\parskip}{0pt}}
\usepackage{bookmark}
\IfFileExists{xurl.sty}{\usepackage{xurl}}{} % add URL line breaks if available
\urlstyle{same}
% fallback for those not using the hyperref driver hyperxmp:
\makeatletter
\@ifundefined{xmpquote}{\newcommand{\xmpquote}[1]{#1}}{}
\makeatother
\hypersetup{
  pdftitle={Single-cell foundation models represent genes in context{,} but the information is co-expression},
  pdfauthor={Ihor Kendiukhov · 2026},
  colorlinks=true,
  linkcolor={NavyBlue},
  filecolor={Maroon},
  citecolor={NavyBlue},
  urlcolor={NavyBlue},
  pdfcreator={LaTeX via pandoc}}

\title{Single-cell foundation models represent genes in context, but the
information is co-expression}
\author{Ihor Kendiukhov · 2026}
\date{2026}

\begin{document}
\maketitle

\subsection{Abstract}\label{abstract}

A single-cell foundation model reads a cell as a sequence of its genes
and produces, at each layer, a contextual representation of every gene
--- the gene's vector given the other genes in that cell. Whether that
representation carries usable information about a gene's role in a
particular cell is unknown, and is the single-cell analogue of
word-sense disambiguation in language models. We give a systematic
answer. Across three architectures (MaxToki, scGPT, STATE) a gene's
representation is reshaped by cellular context in a way that is
gene-specific (cross-cell agreement 0.759 versus 0.001 for a
gene-shuffled null), reproducible across independent cells, not an
artefact of the gene's rank position (96\% survives rank control), and
replicable on an independent developmental dataset. It is causally used:
steering a functional direction into half a cell's genes shifts the
model's own predictions for the other genes (p = 7×10⁻⁴). It is learned
rather than architectural: an untrained model of the same architecture
organises context far less (z = +3.8 versus +21 for the trained model).
It also increases with model size, though this scaling increment is
partly confounded by embedding width. The context modulation is
organised along interpretable functional axes, but against size-matched
co-expression nulls this organisation does not robustly exceed
co-expression --- the strongest axis is only at the p = 0.05 border and
the others are non-significant --- and two independent probes for novel
context-dependent gene function are null. The useful, ceiling-immune
output is a model metric: contextualisation and its functional
organisation differ sharply by tokenisation (rank-based versus
value-binned versus protein-sequence-token) and scale with size. We
report the phenomenon, place a boundary on what it encodes, and give the
metric.

\textbf{Keywords:} single-cell foundation models; contextual gene
representations; mechanistic interpretability; co-expression; gene
embeddings; word-sense disambiguation; model characterisation.

\subsection{1. Introduction}\label{introduction}

A recurring theme in the interpretability of large language models is
that a network trained on a narrow predictive task builds internal
representations of structure it was never taught. Sequence models of
Othello and chess build representations of the board and of latent
variables such as player skill (Li et al. 2023; Nanda et al. 2023;
Karvonen 2024); language models linearly encode space and time along
identifiable directions (Gurnee and Tegmark 2024); and, in the best
cases, these representations are not decorative but \emph{causally used}
--- a claim one tests by intervening on the representation and measuring
the effect on behaviour (Turner et al. 2023; Zou et al. 2023; Park et
al. 2024).

A closely related line concerns \emph{contextual} representations.
Contextual word embeddings assign the same word a different vector in
every sentence, and the geometry of those vectors encodes word sense:
``bank'' near a river and ``bank'' in finance occupy separable regions
(Peters et al. 2018; Ethayarajh 2019; Coenen et al. 2019). A cottage
industry of methodological care surrounds these measurements, because
representation spaces are strongly anisotropic (vectors cluster in a
narrow cone rather than spreading evenly, so any two look similar) and a
handful of ``rogue'' dimensions can dominate cosine similarity
(Ethayarajh 2019; Mu et al. 2018; Cai et al. 2021; Timkey and Schijndel
2021).

Single-cell foundation models (SCFMs) offer a clean biological version
of the same question (Theodoris et al. 2023; Cui et al. 2024; Hao et al.
2024; Yang et al. 2022). They treat each gene as a token and each cell
as a sequence of its expressed genes, and are trained by masked or
autoregressive prediction over expression-ranked genes --- in effect,
learning which genes tend to be active together. Their gene \emph{token}
embeddings (one fixed vector per gene) have been studied; their
\emph{contextual} gene representations (a gene's vector inside a
particular cell) are used as features downstream (Hao et al. 2024; Bian
et al. 2024) but have not been characterised as objects in their own
right. Yet this is exactly where a model could, in principle, represent
that a gene plays a different role in a macrophage than in a T cell ---
the biological analogue of word sense, and, as W. Gilpin has noted, an
open direction that the ``cell as a token'' framing invites but does not
test (W 2025).

We organise the question into two levels. \textbf{Level 1 (a
model-fact):} is a gene's representation reshaped by context in a
structured, gene-specific way, and is that reshaping learned and used?
\textbf{Level 2 (biology):} does the reshaping reveal something
\emph{biologically novel} about the gene's role in each context ---
beyond what co-expression already tells us? We answer Level 1
affirmatively and comprehensively, and Level 2 negatively and with
matched-null rigour. The negative is not a failure to report but the
paper's second contribution: it locates, precisely and with controls,
the ceiling that a growing critical literature has found for SCFMs in
other guises (Kedzierska et al. 2025; Boiarsky et al. 2024;
Ahlmann-Eltze et al. 2025; Kendiukhov 2026e). The paper thus makes three
contributions: (1) the Level-1 model-fact, characterised across
architectures, scale, and a causal test; (2) the Level-2 co-expression
ceiling, established with size-matched nulls; and (3) contextualisation
as a model-characterisation metric that distinguishes tokenisation
schemes (§4.7).

\subsection{2. Related work}\label{related-work}

\textbf{Contextual representations and word sense.} Our methods descend
directly from the NLP interpretability literature on contextualised
embeddings. Ethayarajh (Ethayarajh 2019) showed that representations are
anisotropic in every layer, that same-word self-similarity falls in
upper layers (they become more context-specific), and that a static
vector explains little of a word's contextual variance --- and,
crucially, that every cosine must be corrected against an anisotropy
baseline. Timkey and van Schijndel (Timkey and Schijndel 2021) showed
that a few rogue dimensions dominate similarity and that a scalar
correction is insufficient; Mu et al. (Mu et al. 2018) and Cai et al.
(Cai et al. 2021) give the standard postprocessing and the
local-isotropy picture. Coenen et al. (Coenen et al. 2019) and Garí
Soler and Apidianaki (Soler and Apidianaki 2021) show word senses form
separable clusters and that polysemy level is itself readable. We import
both the measurements and the controls.

\textbf{Emergent, causally-used latent variables.} That a model
represents and \emph{uses} structure it was never given is established
for board games (Li et al. 2023; Nanda et al. 2023; Karvonen 2024) and
for physical coordinates in language models (Gurnee and Tegmark 2024).
We adopt this literature's evidentiary standard --- a probe shows only
correlation, so a causal claim requires an intervention against a
matched control --- and its steering machinery (Turner et al. 2023; Zou
et al. 2023), of which the extraction and amplification of interpretable
feature directions in production models (Templeton et al. 2024) is the
large-scale instance. Our causal result is a steer-here/read-there test
in this mould.

\textbf{Single-cell foundation models and their representations.}
Geneformer (Theodoris et al. 2023), scGPT (Cui et al. 2024),
scFoundation (Hao et al. 2024), scBERT (Yang et al. 2022), UCE (Rosen et
al. 2023) and STATE ({Adduri et al.} 2025) differ in tokenisation
(rank-ordered gene identity, value-binned expression, or
protein-sequence-initialised gene tokens), and the review of Bian et al.
(Bian et al. 2024) defines the ``contextual gene embedding'' our work
dissects. MaxToki ({Ortega et al.} 2026), the primary model here, is a
Llama-architecture, Geneformer-lineage model at 217M and 1B scales.
scGPT reports cell-state-specific gene networks via attention, the
closest prior to a ``context changes the gene'' claim --- but
operationalised as a pairwise network, never as a property of the
representation vector, and without a decomposition or a null. The
\emph{static} gene-embedding geometry of these models has been mapped
--- its spectral structure (Kendiukhov 2026b), the topological and
geometric structure it does and does not learn (Kendiukhov 2026g), and a
decomposition of what a protein language model already supplies versus
what single-cell pretraining adds (Kendiukhov 2026f) --- but the
\emph{contextual} representation, the object here, has not.

\textbf{What SCFMs do not capture.} A critical literature finds SCFMs
often fail to beat simple baselines: on zero-shot tasks (Kedzierska et
al. 2025), under deeper evaluation (Boiarsky et al. 2024), on
perturbation prediction (Ahlmann-Eltze et al. 2025; Bendidi et al.
2024), and across utilities (Liu et al. 2026); and sparse autoencoders
find organised biological knowledge but minimal regulatory logic in
these models (Kendiukhov 2026d), while SCFM attention re-expresses
co-expression rather than unique regulatory signal (Kendiukhov 2026e).
Our co-expression ceiling is the representation-geometry instance of
this pattern, established with a matched null.

\textbf{The biology of the ceiling.} Co-expression is the canonical
guilt-by-association signal (Dam S et al. 2018), and gene-embedding
methods from Gene2vec (Du et al. 2019) to GenePT (YT and J 2023) are
built on it. Genuine context-dependent gene function exists ---
moonlighting proteins (CJ 1999), lineage-determining transcription
factors that select different enhancers with different partners by cell
type (S et al. 2010) --- but much of it is post-transcriptional and so,
we argue, structurally invisible to an expression-only model, which is
why our Level-2 tests are informative negatives.

\subsection{3. Methods}\label{methods}

\textbf{Extraction.} We run cells from Tabula Sapiens (immune, kidney,
lung (Tabula Sapiens Consortium 2022)) through a model and record, for
each gene in a panel, its hidden state at that gene's token position,
averaged over cells of one cell type. The headline MaxToki-217M
extraction covers 12 cell types × 1000 cells, a 6000-gene panel, taps at
layers 0, 4, 8, 11. Four design choices, each fixing a measured failure
of a pilot: (1) \textbf{pairwise contexts}, not the all-context
intersection, which leaves \textasciitilde65 housekeeping genes --- the
``stopwords'' that are \emph{most} context-sensitive (Ethayarajh 2019)
--- whereas pairwise comparison recovers \textasciitilde5,700 genes per
pair; (2) \textbf{cell-level split halves}, each per-gene mean estimated
twice from disjoint cells for an honest noise floor; (3) an
\textbf{occurrence cap} (50 tokens per gene, context, partition),
because a per-gene mean has standard error ∝1/√n and unequal counts
inject context-dependent noise into any interaction term; (4)
\textbf{per-dimension z-scoring} before any cosine, to remove the
rogue-dimension inflation a scalar anisotropy correction misses (Timkey
and Schijndel 2021).

\textbf{Decomposition.} For gene \emph{g} in cell type \emph{c}, write
\emph{v(g,c) = μ + a(g) + b(c) + E(g,c)}: gene main effect \emph{a(g)}
(identity), context main effect \emph{b(c)} (the \textbf{averaging null}
--- attention pools over the cell, so everything drifts together), and
gene-specific interaction \emph{E(g,c)}, the only term where biology can
live.

\textbf{Contextualisation strength (EXCESS).} For a cell-type pair
(c₀,c₁), the \emph{crowd-removed shift} of gene g is δ(g) = {[}v(g,c₁) −
v(g,c₀){]} − mean\_g′{[}v(g′,c₁) − v(g′,c₀){]} --- the gene's own shift
minus the average shift over all genes (the ``crowd'', i.e.~the context
main effect). We compute δ twice, from two disjoint halves of the cells
(δ₀, δ₁), and set EXCESS = mean\_g cos(δ₀(g), δ₁(g)) − mean\_g
cos(δ₀(g), δ₁(perm g)): the same gene's shift agreeing across
independent cells, minus a gene-shuffled control. EXCESS = 0 means pure
averaging (no gene-specific context response); it is averaged over all
cell-type pairs.

\textbf{Functional axes and nulls.} A functional axis u is the unit
vector separating two gene classes (from gene ontology) in the gene main
effect: u ∝ mean(a(g)\textbar class A) − mean(a(g)\textbar class B). Its
\textbf{reproducible interaction power} is: project each gene's two-way
interaction term E(g,c) onto u to get a scalar per (gene, context), then
average the product of the two cell-partitions' scalars over
count-balanced (gene, context) cells --- a covariance that survives only
if the modulation reproduces across independent cells. We compare this
power for u against (i) matched \textbf{random-axis} nulls (random gene
groupings of the same pole sizes); (ii) \textbf{size- and
co-expression-matched} null gene-set pairs --- the decisive control,
since functional gene sets are co-expressed --- where co-expression
coherence is the mean pairwise Pearson correlation of expression across
the extraction cells within the gene set, the null pairs are drawn to
match the functional poles on both size and coherence, and the
functional axis's power is reported as an empirical percentile among
them (skew-robust, because the null power distribution is right-skewed);
and (iii) a \textbf{representation-tightness-matched} null (tightness =
mean cosine among the axis genes' own representations). Full parameters
in §7.

\subsection{4. Results}\label{results}

Sections 4.1--4.3, 4.5 and 4.7--4.8 establish \textbf{Level 1} (the
model-fact: contextualisation is real, not position, replicable, causal,
learned, and scaling); §4.4 and §4.6 test \textbf{Level 2} (novel
biology) and find the co-expression ceiling.

\subsubsection{4.1 Genes are contextualised gene-specifically, and it
peaks in the early-middle
layers}\label{genes-are-contextualised-gene-specifically-and-it-peaks-in-the-early-middle-layers}

The gene-specific context response is large and highly reproducible
(Table 1, Fig. 1). At layer 4 the gene-specific EXCESS is
\textbf{+0.758}: the same gene's crowd-removed shift agrees at cosine
+0.759 across independent cells, while a \emph{different} gene's agrees
at \textbf{+0.001} --- the effect is entirely gene-specific. The context
main effect replicates at +0.996, confirming the measurement works.
Layer 0 (the context-free embedding output) returns \textbf{exactly
0.000}, the sanity check that the pipeline reports nothing when there is
nothing. A scan across layers 0, 4, 8 and 11 places the peak at layer 4,
rising through the early layers and decaying toward the output.

{\def\LTcaptype{none} % do not increment counter
\begin{longtable}[]{@{}
  >{\raggedright\arraybackslash}p{(\linewidth - 8\tabcolsep) * \real{0.2000}}
  >{\raggedright\arraybackslash}p{(\linewidth - 8\tabcolsep) * \real{0.2000}}
  >{\raggedright\arraybackslash}p{(\linewidth - 8\tabcolsep) * \real{0.2000}}
  >{\raggedright\arraybackslash}p{(\linewidth - 8\tabcolsep) * \real{0.2000}}
  >{\raggedright\arraybackslash}p{(\linewidth - 8\tabcolsep) * \real{0.2000}}@{}}
\toprule\noalign{}
\begin{minipage}[b]{\linewidth}\raggedright
Layer
\end{minipage} & \begin{minipage}[b]{\linewidth}\raggedright
EXCESS (same − diff)
\end{minipage} & \begin{minipage}[b]{\linewidth}\raggedright
same-gene
\end{minipage} & \begin{minipage}[b]{\linewidth}\raggedright
diff-gene
\end{minipage} & \begin{minipage}[b]{\linewidth}\raggedright
main-effect replication
\end{minipage} \\
\midrule\noalign{}
\endhead
\bottomrule\noalign{}
\endlastfoot
0 (embedding) & \textbf{+0.000} & +0.000 & +0.000 & --- \\
4 & \textbf{+0.758} & +0.759 & +0.001 & +0.996 \\
8 & +0.665 & +0.666 & +0.001 & +0.996 \\
11 & +0.627 & +0.628 & +0.001 & +0.997 \\
\end{longtable}
}

\emph{Table 1. Gene-specific context response by layer (MaxToki-217M).
66 cell-type pairs; \textasciitilde5,700 genes are shared per pair, of
which a median \textasciitilde1,700 remain after the occurrence-cap
count-balancing (§3); bootstrap CIs ±0.002.}

\begin{figure}
\centering
\includegraphics[width=0.62\linewidth,height=\textheight,keepaspectratio,alt={Figure 1. Gene-specific context response by layer. Same-gene agreement across independent cells (blue) far exceeds the gene-shuffled null (grey); layer 0, the context-free embedding, is exactly zero.}]{figures/ctx_fig1.pdf}
\caption{\textbf{Figure 1.} Gene-specific context response by layer.
Same-gene agreement across independent cells (blue) far exceeds the
gene-shuffled null (grey); layer 0, the context-free embedding, is
exactly zero.}
\end{figure}

\textbf{Metric validity (anisotropy).} These representations are
anisotropic, as expected (Ethayarajh 2019): the mean cosine of random
different-gene pairs is +0.45 to +0.53 at layers 4--11 (raw), and
same-gene self-similarity falls with depth (+0.94 → +0.82 → +0.75) ---
genes become more context-specific in upper layers, the single-cell echo
of Ethayarajh's finding. Our per-dimension z-scoring drives anisotropy
to \textasciitilde0.00 while keeping genes distinguishable (corrected
self-similarity +0.89/+0.65/+0.53), and reduces the top-PC share
(0.21→0.12 at layer 11), so the cosine-based metrics are computed in a
corrected space.

\subsubsection{4.2 It is not the gene's rank
position}\label{it-is-not-the-genes-rank-position}

MaxToki reads a cell as a rank-ordered gene list and encodes token
position, so a gene ranking 5th in one cell type and 500th in another is
an obvious confound. It is not the driver (Fig. 2). The gene-specific
response is essentially identical for rank-stable genes (+0.759) and
rank-moving genes (+0.755), and \textbf{96\%} survives regressing the
projection on per-(gene, context) mean rank (+0.728 residualised at
layer 4). Shift \emph{magnitude} correlates with rank change (ρ ≈ 0.35),
but its \emph{direction} is largely orthogonal to rank.

\begin{figure}
\centering
\includegraphics[width=0.62\linewidth,height=\textheight,keepaspectratio,alt={Figure 2. The effect is not rank position: gene-specific response is identical for rank-stable and rank-moving genes and survives residualising on mean rank.}]{figures/ctx_fig2.pdf}
\caption{\textbf{Figure 2.} The effect is not rank position:
gene-specific response is identical for rank-stable and rank-moving
genes and survives residualising on mean rank.}
\end{figure}

\subsubsection{4.3 It replicates on an independent
dataset}\label{it-replicates-on-an-independent-dataset}

All of the above is on one atlas. On an independent dataset --- Setty
CD34+ bone-marrow hematopoiesis (the same lineage in which a recoverable
manifold was found in scGPT (Kendiukhov 2026a)), with eight
\emph{developmental states} rather than adult cell types as the context
axis --- the phenomenon replicates: EXCESS \textbf{+0.640} (layer 2) and
\textbf{+0.584} (layer 4), gene-shuffled null +0.004, functional-z
\textbf{+9.1/+7.6}. Neither the contextualisation nor its functional
organisation is a Tabula-Sapiens or adult-cell-type artefact.

\subsubsection{4.4 The modulation is functionally organised --- but the
organisation is
co-expression}\label{the-modulation-is-functionally-organised-but-the-organisation-is-co-expression}

Context modulates genes along interpretable functional axes far more
than along random directions. Reported as raw z (rank-controlled in
parentheses), the nuclear/transcriptional vs.~surface/secreted axis
reaches z = \textbf{+21.9} (+20.2); mitochondrial vs.~cytoskeletal +10.3
(+8.3); transcription vs.~transport +7.1 (+6.3); the abundance-neutral
but weakly-separable cell-cycle-vs-differentiation axis is null, +0.7
(−0.7) at AUC 0.54. The nuclear/surface axis is biologically valid ---
cell types order sensibly, lymphocytes at the nuclear pole (CD8 T +4.1,
CD4 T +3.7, B +3.4) and secretory/surface cells at the other (alveolar
type-2 −5.3, type-1 −3.7, macrophage −3.2).

The decisive control asks whether this exceeds \textbf{co-expression}
(Fig. 3). Functional gene sets are co-regulated modules, so their genes
co-vary across cells regardless of any ``understanding'' of function.
The right null must match the functional poles on \textbf{both} their
size and their co-expression: comparing a \textasciitilde1,880-gene
functional axis against small random modules is invalid, because smaller
sets reach higher coherence and different power. We therefore build null
gene-set pairs of the \emph{same sizes} as the functional poles, matched
on co-expression coherence, and --- because the null power distribution
is heavily right-skewed --- report the functional axis's power as an
empirical percentile among them. The nuclear/surface axis's power (5.32)
sits at the \textbf{p = 0.05} border of size-matched co-expression
modules (mean 2.22); the mitochondrial/cytoskeletal axis at p = 0.16 and
the transcription/transport axis at p = 0.10 are not significant. So the
functional organisation does \emph{not} robustly exceed co-expression
--- the headline axis is at best marginal, the others indistinguishable.
It does clear a representation-space \textbf{tightness}-matched null
(mean cosine among the axis genes' own representations) and, by
construction, the random-axis null (p \textless{} 0.001), so it is not
merely a compact-gene-set artefact --- but co-expression is sufficient.
This is the geometry instance of the field-wide pattern that SCFMs
re-express co-expression (Kendiukhov 2026e; Kedzierska et al. 2025).

\begin{figure}
\centering
\includegraphics[width=0.66\linewidth,height=\textheight,keepaspectratio,alt={Figure 3. The crux control. Random gene modules (grey) trace context-modulation power against co-expression coherence; the functional axes (red) sit on that curve, so the functional organisation is co-expression.}]{figures/ctx_fig3.pdf}
\caption{\textbf{Figure 3.} The crux control. Random gene modules (grey)
trace context-modulation power against co-expression coherence; the
functional axes (red) sit on that curve, so the functional organisation
is co-expression.}
\end{figure}

\subsubsection{4.5 The functional-context direction is causally
used}\label{the-functional-context-direction-is-causally-used}

Level 1 is a genuine model-fact, and it is \emph{used}, not decorative
(Fig. 4). In a steer-here/read-there test we add a functional direction
\emph{u} to a random half of a cell's genes and read the model's own
output logits at the \emph{other}, untouched half. Nuclear-gene logits
rise toward +\emph{u} (specific effect \textbf{+0.160}) and fall toward
−\emph{u} (\textbf{−0.112}), a signed swing of \textbf{+0.272} versus
\textbf{−0.015} for a norm-matched random push, dose-dependently across
three strengths and beyond the random control on 24 of 30 cells
(\emph{p} = 7×10⁻⁴). The functional-context direction propagates through
attention and shifts the model's predictions.

The effect is not specific to one axis or layer. It replicates for the
same nuclear/surface direction injected at a deeper layer (layer 8:
swing +0.566, 30/30 cells, \emph{p} = 9×10⁻¹⁰, stronger than at layer 4)
and for an independent functional axis (mitochondrial vs.~cytoskeletal
at layer 4: swing +0.321, \emph{p} = 0.02). It is null for the
weakest-separated axis we tested causally (transcription vs.~transport,
validity AUC 0.66: swing −0.226, \emph{p} = 1.0 --- in fact a reversal),
consistent with a channel that exists for well-formed functional
directions but not for every ontology contrast one can write down.
Across the four causal configurations the p-values are raw; the two
strongest (nuclear/surface at layers 4 and 8) survive Bonferroni
correction for the four tests, while the mitochondrial axis (p = 0.02)
becomes marginal.

\begin{figure}
\centering
\includegraphics[width=0.62\linewidth,height=\textheight,keepaspectratio,alt={Figure 4. Causal use. Steering a functional direction into half a cell's genes shifts the model's logits at the other genes with the sign of the push, dose-dependently, far beyond a norm-matched random push.}]{figures/ctx_fig4.pdf}
\caption{\textbf{Figure 4.} Causal use. Steering a functional direction
into half a cell's genes shifts the model's logits at the \emph{other}
genes with the sign of the push, dose-dependently, far beyond a
norm-matched random push.}
\end{figure}

\subsubsection{4.6 No novel context-dependent gene function (the
co-expression
ceiling)}\label{no-novel-context-dependent-gene-function-the-co-expression-ceiling}

Two independent attempts to read \emph{context-appropriate or novel}
gene function beyond co-expression are null. (1) \textbf{Directional
congruence} --- does a gene move toward its own functional pole
specifically in contexts sharing that character (a cross-validated
non-additivity test)? Inside the co-expression null (empirical p ≈ 0.2);
no signal. (2) \textbf{Placement near curated targets} --- are
transcription factors positioned near their ChIP-seq / literature
targets, context-appropriately (21 TFs)? Static excess +0.023 (\emph{p}
= 0.19), activity-modulation ρ = +0.080 (\emph{p} = 0.33); no signal,
even with real target ground truth. A third attempt (curated
context-switching TFs) is uninformative rather than null: only one such
TF (RUNX1) is well-sampled in the panel, because context- switching TFs
are typically low-expressed --- itself an instance of the measurement
wall. The boundary is consistent: an expression-only SCFM contextualises
genes richly, but the information is co-expression, and does not resolve
a gene's changed role in a new cell beyond it (the ceiling is quantified
in Fig. 3).

\subsubsection{4.7 Contextualisation as a model metric: architecture,
scale, and learned
vs.~architectural}\label{contextualisation-as-a-model-metric-architecture-scale-and-learned-vs.-architectural}

The direction immune to the co-expression ceiling is to treat the
phenomenon as a \textbf{model-quality metric} (Table 2, Fig. 5). Two
findings. (1) \textbf{Contextualisation is universal and scales.} Every
architecture strongly contextualises genes (EXCESS 0.70--0.88,
main-effect replication \textgreater{} 0.93); on matched data (600
cells, identical settings) it increases with size, MaxToki-217M +0.740 →
1B +0.881, in line with the loss-scaling behaviour reported for these
masked-reconstruction models (Kendiukhov 2026c). (2)
\textbf{Architecture determines functional organisation.} Rank-based
MaxToki organises context modulation along functional axes an order of
magnitude more (z = 16--22) than value-binned scGPT (+2.7) or
protein-sequence-initialised STATE (+1.3). STATE is telling: it moves
genes strongly in context (EXCESS +0.835) but in no functionally
coherent direction --- its ESM2-initialised gene tokens carry sequence
identity, not co-expression. The gap is real, not low power (pole sizes
366--1,390).

\begin{enumerate}
\def\labelenumi{(\arabic{enumi})}
\setcounter{enumi}{2}
\tightlist
\item
  \textbf{The phenomenon is learned, not architectural.} An untrained,
  random-init MaxToki of identical architecture and vocabulary --- the
  decisive control, since attention mixes tokens regardless of training
  --- contextualises genes far less (EXCESS \textbf{+0.274}
  vs.~+0.74--0.76 trained) and organises that modulation functionally
  far less (functional-z \textbf{+3.8} vs.~+21.2). This separates two
  things cleanly. Contextualisation has a real architectural
  \emph{floor} (+0.274: fixed random embeddings plus deterministic
  attention mixing give a reproducible, gene-specific shift on their
  own), but roughly two-thirds of the trained effect is learned. The
  functional organisation is \emph{mostly} learned: the architecture
  alone produces a significant amount (z = +3.8), which training then
  amplifies \textasciitilde5.6× (to +21.2) --- so roughly 18\% of the
  functional organisation is architectural and 82\% is built by
  training, aligning context modulation with functional axes, and hence
  with co-expression.
\end{enumerate}

{\def\LTcaptype{none} % do not increment counter
\begin{longtable}[]{@{}
  >{\raggedright\arraybackslash}p{(\linewidth - 12\tabcolsep) * \real{0.1429}}
  >{\raggedright\arraybackslash}p{(\linewidth - 12\tabcolsep) * \real{0.1429}}
  >{\raggedright\arraybackslash}p{(\linewidth - 12\tabcolsep) * \real{0.1429}}
  >{\raggedright\arraybackslash}p{(\linewidth - 12\tabcolsep) * \real{0.1429}}
  >{\raggedright\arraybackslash}p{(\linewidth - 12\tabcolsep) * \real{0.1429}}
  >{\raggedright\arraybackslash}p{(\linewidth - 12\tabcolsep) * \real{0.1429}}
  >{\raggedright\arraybackslash}p{(\linewidth - 12\tabcolsep) * \real{0.1429}}@{}}
\toprule\noalign{}
\begin{minipage}[b]{\linewidth}\raggedright
Model
\end{minipage} & \begin{minipage}[b]{\linewidth}\raggedright
tokenisation
\end{minipage} & \begin{minipage}[b]{\linewidth}\raggedright
contexts
\end{minipage} & \begin{minipage}[b]{\linewidth}\raggedright
genes
\end{minipage} & \begin{minipage}[b]{\linewidth}\raggedright
dim
\end{minipage} & \begin{minipage}[b]{\linewidth}\raggedright
EXCESS
\end{minipage} & \begin{minipage}[b]{\linewidth}\raggedright
functional-z
\end{minipage} \\
\midrule\noalign{}
\endhead
\bottomrule\noalign{}
\endlastfoot
scGPT & value-binned & 8 & 5000 & 512 & +0.705 & +2.7 \\
STATE-SE (state-embedding) & ESM2 gene tokens & 9 & 1840 & 2048 & +0.835
& +1.3 \\
MaxToki-217M & rank-based & 12 & 5000 & 1232 & +0.740 & +21.2 \\
MaxToki-1B & rank-based & 12 & 5000 & 2304 & +0.881 & +16.4 \\
MaxToki-217M (random init) & untrained control & 12 & 6000 & 1232 &
+0.274 & +3.8 \\
\end{longtable}
}

\emph{Table 2. Cross-model comparison at layer 4. EXCESS =
contextualisation strength; functional-z = organisation along the
nuclear/surface axis vs.~a random-axis null. MaxToki-217M/1B use matched
600-cell settings for a fair scaling comparison; the random-init row is
the learned-vs-architectural control (§4.7); its 6000-gene panel vs. the
5000 of the trained rows does not affect EXCESS, which is a within-panel
cosine statistic.}

\begin{figure}
\centering
\includegraphics[width=0.92\linewidth,height=\textheight,keepaspectratio,alt={Figure 5. Contextualisation as a model metric across architecture, scale, and the random-weights control.}]{figures/ctx_fig5.pdf}
\caption{\textbf{Figure 5.} Contextualisation as a model metric across
architecture, scale, and the random-weights control.}
\end{figure}

\subsubsection{4.8 Context aids prediction, but not through the
most-contextualised
genes}\label{context-aids-prediction-but-not-through-the-most-contextualised-genes}

If contextualisation served the autoregressive objective directly, the
genes the model contextualises most should be the genes whose next-gene
prediction most benefits from real context. It does not work that way.
Using a chimeric-context intervention --- the model's log-probability of
a target gene given its cell's true rank-ordered predecessors versus a
random other cell's --- real context helps prediction enormously (mean
benefit \textbf{+7.7 nats}, a nat being a natural-log unit of
log-probability; i.e.~the true gene is roughly two-thousand-fold more
probable under real than random context; 536 target genes). But the
per-gene strength of that benefit does \textbf{not} track how much a
gene is contextualised: the partial Spearman correlation, controlling
for gene frequency, is \textbf{−0.087} (95\% CI {[}−0.168, −0.004{]})
--- at best faintly negative, and with per-gene estimates noisy on both
sides, best read as \emph{no positive link} rather than a firm negative
one. Context is central to the model's objective, and contextualisation
is real, but the two are not aligned per gene: contextualisation is not
simply the model encoding its per-gene predictive reliance on context.

\subsection{5. Discussion}\label{discussion}

The result is a characterisation of \emph{how} SCFMs represent genes in
context, and a proof of \emph{what} that representation carries.
Positively: SCFMs reshape gene representations gene-specifically,
reproducibly, causally, more with scale, and --- for rank-based
tokenisation --- along co-expression-aligned functional axes. This is,
to our knowledge, the first measurement for single-cell models of the
``how contextual are contextualised representations'' question
(Ethayarajh 2019), and the architecture/scale comparison is a usable
interpretability metric that discriminates tokenisation schemes a
downstream benchmark would not.

Negatively, and equally a contribution: the modulation does not carry a
gene's context-specific function beyond co-expression. The exciting
version --- reading a gene's changed role in a new cell to generate
hypotheses --- is not reachable for an expression-only model by either
route we tried, each with a matched-null control. Why the ceiling exists
is mechanistically coherent: the model's only inputs are which genes are
co-expressed, so the richest thing its contextual geometry can encode
about a gene is its co-expression neighbourhood; genuinely novel
context-dependent function is frequently post-transcriptional (CJ 1999;
S et al. 2010) and leaves no expression signature to latch onto. The
finding therefore predicts where headroom is: multi-omic or
sequence-aware models, not larger expression-only ones, are what a
Level-2 result would require.

\subsection{6. Limitations}\label{limitations}

Cross-architecture EXCESS is confounded by occurrence cap (20 for
scGPT/STATE vs.~50 for MaxToki) and dimensionality (higher-dimensional
models show higher EXCESS), so the 217M→1B increment reflects size
\emph{and} width, and even the cap-matched scaling pair still differs in
width (dim 1232 vs.~2304) --- the width-robust comparison is the
functional-z architecture gap, which has ample power. The functional-z
gap itself is measured on non-identical context sets across models (12
vs.~9 vs.~8 cell types); we did not re-run it on a common context set,
so a residual context-composition confound remains. The causal test
(§4.5) compares the functional direction only against a norm-matched
\emph{random} push, not against co-expression-coherence-matched
directions, so it shows the direction is used, not that its causal use
exceeds co-expression. STATE is evaluated on immune cells only in this
study, so its cross-architecture point rests on a narrower panel than
the others. The co-expression ceiling is a null result: the headline
axis sits at the p = 0.05 border of the size-matched co-expression null
and the other axes are non-significant, which shows the organisation
does not clearly exceed co-expression, not that the two are provably
identical; the empirical p depends on the null-module construction. The
Level-2 null is evidence of a ceiling, not proof of absence --- some
context-dependent function is structurally invisible to an
expression-only model, so those tests are weak by construction
(pre-registered). Functional axes are built from one ontology file; the
verdict is robust across three axes but specific z-scores depend on pole
definitions.

\subsection{7. Conclusions}\label{conclusions}

Single-cell foundation models build contextual gene representations that
are a genuine object of study: a gene's vector is reshaped by its
cellular context in a gene-specific, reproducible, learned, and causally
used way that grows with model scale and is shaped by tokenisation. This
is the first systematic characterisation of the phenomenon for
single-cell models, and it yields a usable model metric ---
contextualisation strength and its functional organisation --- that
distinguishes architectures a downstream benchmark would not. At the
same time, matched-null controls place a firm boundary on what that
contextualisation carries: it does not robustly exceed co-expression,
and two independent probes for novel context-dependent gene function are
null. The practical implication is that reading a gene's changed role in
a new cell type --- the application that would make these
representations a hypothesis-generation tool --- is not reachable with
an expression-only model, and will require multi-omic or sequence-aware
models rather than merely larger expression-trained ones.

\subsection{8. Methods details and
reproducibility}\label{methods-details-and-reproducibility}

All analyses are deterministic (seed 0) and run on cached extractions.
Extraction: \texttt{ctx\_extract\_maxtoki.py} (+ \texttt{\_scgpt},
\texttt{\_state}, \texttt{\_devel}, \texttt{\_random}). Analyses:
\texttt{ctx\_polysemy.py} (§4.1), \texttt{ctx\_anisotropy.py} (§4.1),
\texttt{ctx\_layer\_curve.py} (§4.1),
\texttt{ctx\_position\_confound.py} (§4.2), \texttt{ctx\_independent.py}
(§4.3), \texttt{ctx\_functional\_axes.py},
\texttt{ctx\_coexpr\_null\_v2.py}, \texttt{ctx\_tightness\_null.py}
(§4.4), \texttt{ctx\_causal.py} across axes/layers (§4.5),
\texttt{ctx\_directional\_probe.py} / \texttt{ctx\_curated\_targets.py}
(§4.6), \texttt{ctx\_cross\_model.py} (§4.7),
\texttt{ctx\_prediction\_link.py} (§4.8). Metric definitions: EXCESS as
in §3; reproducible interaction power = mean over count-balanced (gene,
context) of the product of the two partitions' two-way interaction terms
along a unit axis; functional-z, its z-score against 200--300
random-partition axes; co-expression coherence = mean within-set
expression correlation. The co-expression null
(\texttt{ctx\_coexpr\_null\_v2.py}) draws null gene-set pairs of the
\emph{same sizes} as the functional poles, matched on coherence, and
reports the functional axis's power as an empirical percentile among
them (skew-robust). Result artefacts are the corresponding
\texttt{results/ctx\_*.json}.

\subsection{CRediT authorship contribution
statement}\label{credit-authorship-contribution-statement}

\textbf{Ihor Kendiukhov:} Conceptualization, Methodology, Software,
Formal analysis, Investigation, Data curation, Writing -- original
draft, Writing -- review \& editing, Visualization.

\subsection{Declaration of competing
interest}\label{declaration-of-competing-interest}

The author declares that he has no known competing financial interests
or personal relationships that could have appeared to influence the work
reported in this paper.

\subsection{Data and code
availability}\label{data-and-code-availability}

All analysis code is openly available at
\url{https://github.com/Biodyn-AI/gene-context}, including every
analysis script, the intermediate result artefacts
(\texttt{results/ctx\_*.json}), the figure-generation code, and this
manuscript with its bibliography. The study uses only publicly available
models and data: the MaxToki-217M/1B checkpoints (Hugging Face,
\texttt{theodoris-lab/MaxToki}); scGPT, Geneformer and STATE (public
releases); the Tabula Sapiens atlas (Tabula Sapiens Consortium 2022) (CZ
CELLxGENE); the Setty CD34+ bone-marrow dataset; and Gene Ontology
annotations. The per-gene contextual representation tensors extracted
from these models (the intermediate \texttt{.npz} files, several
gigabytes) are not deposited because of their size, but are fully
regenerable from the public inputs with the provided extraction scripts,
and are available from the author on reasonable request. No new
experimental data were generated.

\subsection{Funding}\label{funding}

This research did not receive any specific grant from funding agencies
in the public, commercial, or not-for-profit sectors.

\subsection{References}\label{references}

\section*{References}\label{bibliography}
\addcontentsline{toc}{section}{References}

\protect\phantomsection\label{refs}
\begin{CSLReferences}{1}{1}
\bibitem[\citeproctext]{ref-state2025}
{Adduri, Abhinav K., Dhruv Gautam, Beatrice Bevilacqua, et al.} 2025.
{``Predicting Cellular Responses to Perturbation Across Diverse Contexts
with State.''} \emph{bioRxiv (Preprint), Arc Institute}, ahead of print.
\url{https://doi.org/10.1101/2025.06.26.661135}.

\bibitem[\citeproctext]{ref-ahlmanneltze2025perturbation}
Ahlmann-Eltze, Constantin, Wolfgang Huber, and Simon Anders. 2025.
{``Deep-Learning-Based Gene Perturbation Effect Prediction Does Not yet
Outperform Simple Linear Baselines.''} \emph{Nature Methods
22(8):1657-1661 (Preprint bioRxiv 10.1101/2024.09.16.613342; PubMed
40759747)}, ahead of print.
\url{https://doi.org/10.1038/s41592-025-02772-6}.

\bibitem[\citeproctext]{ref-bendidi2024pca}
Bendidi, Ihab, Shawn Whitfield, Kian Kenyon-Dean, et al. 2024.
{``Benchmarking Transcriptomics Foundation Models for Perturbation
Analysis: One PCA Still Rules Them All.''} \emph{arXiv:2410.13956
(NeurIPS 2024 Workshop Paper)}. \url{https://arxiv.org/abs/2410.13956}.

\bibitem[\citeproctext]{ref-bian2024}
Bian, Haiyang, Yixin Chen, Erpai Luo, et al. 2024. {``General-Purpose
Pre-Trained Large Cellular Models for Single-Cell Transcriptomics.''}
\emph{National Science Review 11(11):nwae340}, ahead of print.
\url{https://doi.org/10.1093/nsr/nwae340}.

\bibitem[\citeproctext]{ref-boiarsky2024deeper}
Boiarsky, Rebecca, Nalini M. Singh, Alejandro Buendia, Ava P. Amini, Gad
Getz, and David Sontag. 2024. {``Deeper Evaluation of a Single-Cell
Foundation Model.''} \emph{Nature Machine Intelligence 6(12):1443-1446
(Preprint: 'A Deep Dive into Single-Cell RNA Sequencing Foundation
Models', bioRxiv 10.1101/2023.10.19.563100)}, ahead of print.
\url{https://doi.org/10.1038/s42256-024-00949-w}.

\bibitem[\citeproctext]{ref-cai2021}
Cai, Xingyu, Jiaji Huang, Yuchen Bian, and Kenneth Church. 2021.
{``Isotropy in the Contextual Embedding Space: Clusters and
Manifolds.''} \emph{ICLR 2021 (International Conference on Learning
Representations)}. \url{https://openreview.net/forum?id=xYGNO86OWDH}.

\bibitem[\citeproctext]{ref-Jeffery1999Moonlighting}
CJ, Jeffery. 1999. {``Moonlighting Proteins.''} \emph{Trends in
Biochemical Sciences 24(1): 8-11}, ahead of print.
\url{https://doi.org/10.1016/S0968-0004(98)01335-8}.

\bibitem[\citeproctext]{ref-coenen2019}
Coenen, Andy, Emily Reif, Ann Yuan, et al. 2019. {``Visualizing and
Measuring the Geometry of BERT.''} \emph{NeurIPS 2019 (Advances in
Neural Information Processing Systems 32)}.
\url{https://arxiv.org/abs/1906.02715}.

\bibitem[\citeproctext]{ref-scgpt2024}
Cui, Haotian, Chloe Wang, Hassaan Maan, et al. 2024. {``scGPT: Toward
Building a Foundation Model for Single-Cell Multi-Omics Using Generative
AI.''} \emph{Nature Methods 21(8):1470-1480}, ahead of print.
\url{https://doi.org/10.1038/s41592-024-02201-0}.

\bibitem[\citeproctext]{ref-vanDam2018Coexpression}
Dam S, van, Võsa U, van der Graaf A, Franke L, and de Magalhães JP.
2018. {``Gene Co-Expression Analysis for Functional Classification and
Gene-Disease Predictions.''} \emph{Briefings in Bioinformatics 19(4):
575-592}, ahead of print. \url{https://doi.org/10.1093/bib/bbw139)}.

\bibitem[\citeproctext]{ref-Du2019Gene2vec}
Du, Jingcheng, Peilin Jia, Yulin Dai, Cui Tao, Zhongming Zhao, and Degui
Zhi. 2019. {``Gene2vec: Distributed Representation of Genes Based on
Co-Expression.''} \emph{BMC Genomics 20 (Suppl 1): 82}, ahead of print.
\url{https://doi.org/10.1186/s12864-018-5370-x}.

\bibitem[\citeproctext]{ref-ethayarajh2019}
Ethayarajh, Kawin. 2019. {``How Contextual Are Contextualized Word
Representations? Comparing the Geometry of BERT, ELMo, and GPT-2
Embeddings.''} \emph{EMNLP-IJCNLP 2019 (Proceedings of the 2019
Conference on Empirical Methods in Natural Language Processing)}, ahead
of print. \url{https://doi.org/10.18653/v1/D19-1006}.

\bibitem[\citeproctext]{ref-gurnee2024language}
Gurnee, Wes, and Max Tegmark. 2024. {``Language Models Represent Space
and Time.''} \emph{ICLR 2024 (arXiv:2310.02207)}.
\url{https://arxiv.org/abs/2310.02207}.

\bibitem[\citeproctext]{ref-scfoundation2024}
Hao, Minsheng, Jing Gong, Xin Zeng, et al. 2024. {``Large-Scale
Foundation Model on Single-Cell Transcriptomics.''} \emph{Nature Methods
21(8):1481-1491}, ahead of print.
\url{https://doi.org/10.1038/s41592-024-02305-7}.

\bibitem[\citeproctext]{ref-karvonen2024emergent}
Karvonen, Adam. 2024. {``Emergent World Models and Latent Variable
Estimation in Chess-Playing Language Models.''} \emph{arXiv Preprint
arXiv:2403.15498 (2024)}. \url{https://arxiv.org/abs/2403.15498}.

\bibitem[\citeproctext]{ref-kedzierska2025zeroshot}
Kedzierska, Kasia Z., Lorin Crawford, Ava P. Amini, and Alex X. Lu.
2025. {``Zero-Shot Evaluation Reveals Limitations of Single-Cell
Foundation Models.''} \emph{Genome Biology 26(1):101}, ahead of print.
\url{https://doi.org/10.1186/s13059-025-03574-x}.

\bibitem[\citeproctext]{ref-kend2026manifold}
Kendiukhov, Ihor. 2026a. {``Discovery of a Hematopoietic Manifold in
{scGPT} Yields a Method for Extracting Performant Algorithms from
Biological Foundation Model Internals.''} \emph{arXiv Preprint
arXiv:2603.10261}. \url{https://arxiv.org/abs/2603.10261}.

\bibitem[\citeproctext]{ref-kend2026spectral}
Kendiukhov, Ihor. 2026b. {``Multi-Dimensional Spectral Geometry of
Biological Knowledge in Single-Cell Transformer Representations.''}
\emph{arXiv Preprint arXiv:2602.22247}.
\url{https://arxiv.org/abs/2602.22247}.

\bibitem[\citeproctext]{ref-kend2026scaling}
Kendiukhov, Ihor. 2026c. {``Scaling Laws for Masked-Reconstruction
Transformers on Single-Cell Transcriptomics.''} \emph{Transactions on
Machine Learning Research}.

\bibitem[\citeproctext]{ref-kend2026sae}
Kendiukhov, Ihor. 2026d. {``Sparse Autoencoders Reveal Organized
Biological Knowledge but Minimal Regulatory Logic in Single-Cell
Foundation Models: A Comparative Atlas of {Geneformer} and {scGPT}.''}
\emph{arXiv Preprint arXiv:2603.02952}.
\url{https://arxiv.org/abs/2603.02952}.

\bibitem[\citeproctext]{ref-kendiukhov2026attention}
Kendiukhov, Ihor. 2026e. {``Systematic Evaluation of Single-Cell
Foundation Model Interpretability Reveals Attention Captures
Co-Expression Rather Than Unique Regulatory Signal.''}
\emph{arXiv:2602.17532}. \url{https://arxiv.org/abs/2602.17532}.

\bibitem[\citeproctext]{ref-kend2026twoaxes}
Kendiukhov, Ihor. 2026f. {``Two Axes in the Gene-Embedding Space of
Single-Cell Foundation Models: What a Protein Language Model Already
Knows, and What Requires Single-Cell Pretraining.''} \emph{SSRN Preprint
6577900}.

\bibitem[\citeproctext]{ref-kend2026topo141}
Kendiukhov, Ihor. 2026g. {``What Topological and Geometric Structure Do
Biological Foundation Models Learn? {Evidence} from 141 Hypotheses.''}
\emph{PLOS ONE} 21 (7): e0344826.
\url{https://doi.org/10.1371/journal.pone.0344826}.

\bibitem[\citeproctext]{ref-li2023emergent}
Li, Kenneth, Aspen K. Hopkins, David Bau, Fernanda Viégas, Hanspeter
Pfister, and Martin Wattenberg. 2023. {``Emergent World Representations:
Exploring a Sequence Model Trained on a Synthetic Task.''} \emph{ICLR
2023 (arXiv:2210.13382)}. \url{https://arxiv.org/abs/2210.13382}.

\bibitem[\citeproctext]{ref-liu2026sceval}
Liu, Tianyu, Kexing Li, Yuge Wang, Hongyu Li, and Hongyu Zhao. 2026.
{``Evaluating the Utilities of Foundation Models in Single-Cell Data
Analysis.''} \emph{Advanced Science 13(27):e14490 (Preprint bioRxiv
10.1101/2023.09.08.555192)}, ahead of print.
\url{https://doi.org/10.1002/advs.202514490}.

\bibitem[\citeproctext]{ref-muviswanath2018}
Mu, Jiaqi, Suma Bhat, and Pramod Viswanath. 2018. {``All-but-the-Top:
Simple and Effective Postprocessing for Word Representations.''}
\emph{ICLR 2018 (International Conference on Learning Representations)}.
\url{https://arxiv.org/abs/1702.01417}.

\bibitem[\citeproctext]{ref-nanda2023emergent}
Nanda, Neel, Andrew Lee, and Martin Wattenberg. 2023. {``Emergent Linear
Representations in World Models of Self-Supervised Sequence Models.''}
\emph{BlackboxNLP Workshop at EMNLP 2023 (arXiv:2309.00941)}.
\url{https://arxiv.org/abs/2309.00941}.

\bibitem[\citeproctext]{ref-maxtoki2026}
{Ortega, Javier Gomez, Rangarajan D. Nadadur, Akira Kunitomi, et al.}
2026. {``Temporal AI Model Predicts Drivers of Cell State Trajectories
Across Human Aging.''} \emph{bioRxiv (Preprint), 2026.03.30.715396},
ahead of print. \url{https://doi.org/10.64898/2026.03.30.715396}.

\bibitem[\citeproctext]{ref-park2024linear}
Park, Kiho, Yo Joong Choe, and Victor Veitch. 2024. {``The Linear
Representation Hypothesis and the Geometry of Large Language Models.''}
\emph{ICML 2024, PMLR V235, Pp. 39643--39666 (arXiv:2311.03658)}.
\url{https://arxiv.org/abs/2311.03658}.

\bibitem[\citeproctext]{ref-peters2018}
Peters, Matthew E., Mark Neumann, Mohit Iyyer, et al. 2018. {``Deep
Contextualized Word Representations.''} \emph{NAACL-HLT 2018
(Proceedings of the 2018 Conference of the North American Chapter of the
ACL); Best Paper}, ahead of print.
\url{https://doi.org/10.18653/v1/N18-1202}.

\bibitem[\citeproctext]{ref-uce2023}
Rosen, Yanay, Yusuf Roohani, Ayush Agrawal, et al. 2023. {``Universal
Cell Embeddings: A Foundation Model for Cell Biology.''} \emph{bioRxiv
(Preprint); Later Published in Nature (2026)}, ahead of print.
\url{https://doi.org/10.1101/2023.11.28.568918}.

\bibitem[\citeproctext]{ref-Heinz2010LDTFs}
S, Heinz, Benner C, Spann N, et al. 2010. {``Simple Combinations of
Lineage-Determining Transcription Factors Prime Cis-Regulatory Elements
Required for Macrophage and b Cell Identities.''} \emph{Molecular Cell
38(4): 576-589}, ahead of print.
\url{https://doi.org/10.1016/j.molcel.2010.05.004}.

\bibitem[\citeproctext]{ref-garisoler2021}
Soler, Aina Garí, and Marianna Apidianaki. 2021. {``Let's Play
Mono-Poly: BERT Can Reveal Words' Polysemy Level and Partitionability
into Senses.''} \emph{TACL 2021 (Transactions of the Association for
Computational Linguistics, Vol. 9)}, ahead of print.
\url{https://doi.org/10.1162/tacl_a_00400}.

\bibitem[\citeproctext]{ref-tabulasapiens2022}
Tabula Sapiens Consortium. 2022. {``The {Tabula Sapiens}: A
Multiple-Organ, Single-Cell Transcriptomic Atlas of Humans.''}
\emph{Science} 376 (6594): eabl4896.
\url{https://doi.org/10.1126/science.abl4896}.

\bibitem[\citeproctext]{ref-templeton2024scaling}
Templeton, Adly, Tom Conerly, Jonathan Marcus, et al. 2024. {``Scaling
Monosemanticity: Extracting Interpretable Features from Claude 3
Sonnet.''} \emph{Transformer Circuits Thread (2024),
Transformer-Circuits.pub}.
\url{https://transformer-circuits.pub/2024/scaling-monosemanticity/}.

\bibitem[\citeproctext]{ref-geneformer2023}
Theodoris, Christina V., Ling Xiao, Anant Chopra, et al. 2023.
{``Transfer Learning Enables Predictions in Network Biology.''}
\emph{Nature 618(7965):616-624}, ahead of print.
\url{https://doi.org/10.1038/s41586-023-06139-9}.

\bibitem[\citeproctext]{ref-timkey2021}
Timkey, William, and Marten van Schijndel. 2021. {``All Bark and No
Bite: Rogue Dimensions in Transformer Language Models Obscure
Representational Quality.''} \emph{EMNLP 2021 (Proceedings of the 2021
Conference on Empirical Methods in Natural Language Processing)}, ahead
of print. \url{https://doi.org/10.18653/v1/2021.emnlp-main.372}.

\bibitem[\citeproctext]{ref-turner2023activation}
Turner, Alexander Matt, Lisa Thiergart, Gavin Leech, et al. 2023.
{``Activation Addition: Steering Language Models Without Optimization
(ActAdd); Later Revised on arXiv as 'Steering Language Models with
Activation Engineering'.''} \emph{arXiv Preprint arXiv:2308.10248
(2023)}. \url{https://arxiv.org/abs/2308.10248}.

\bibitem[\citeproctext]{ref-Gilpin2025CellToken}
W, Gilpin. 2025. {``The Cell as a Token: High-Dimensional Geometry in
Language Models and Cell Embeddings.''} \emph{Bioinformatics 41(11):
Btaf595}, ahead of print.
\url{https://doi.org/10.1093/bioinformatics/btaf595}.

\bibitem[\citeproctext]{ref-scbert2022}
Yang, Fan, Wenchuan Wang, Fang Wang, et al. 2022. {``scBERT as a
Large-Scale Pretrained Deep Language Model for Cell Type Annotation of
Single-Cell RNA-Seq Data.''} \emph{Nature Machine Intelligence
4(10):852-866}, ahead of print.
\url{https://doi.org/10.1038/s42256-022-00534-z}.

\bibitem[\citeproctext]{ref-Chen2023GenePT}
YT, Chen, and Zou J. 2023. {``GenePT: A Simple but Effective Foundation
Model for Genes and Cells Built from ChatGPT (Preprint); Journal
Version: Simple and Effective Embedding Model for Single-Cell Biology
Built from ChatGPT.''} \emph{bioRxiv 2023.10.16.562533; Journal Version
in Nature Biomedical Engineering (2025)}, ahead of print.
\url{https://doi.org/10.1101/2023.10.16.562533}.

\bibitem[\citeproctext]{ref-zou2023representation}
Zou, Andy, Long Phan, Sarah Chen, et al. 2023. {``Representation
Engineering: A Top-down Approach to AI Transparency.''} \emph{arXiv
Preprint arXiv:2310.01405 (2023)}.
\url{https://arxiv.org/abs/2310.01405}.

\end{CSLReferences}

\end{document}
